% !TEX program = lualatex
% Compile from hackathon/: lualatex participant_warmup.tex
\documentclass[10pt,a4paper]{article}
\usepackage[english]{babel}
\usepackage{tensygrid-hackathon}
\usepackage{amsmath,amssymb,booktabs,tabularx,enumitem}
\usepackage{titlesec,fancyhdr,lastpage}
\usepackage[colorlinks=true,urlcolor=tgpetrol,linkcolor=tgblue]{hyperref}
\pagestyle{fancy}
\fancyhf{}
\lhead{\small\textcolor{tgpetrol}{\textbf{TenSyGrid $\times$ VeraGrid}}\enspace|\enspace Warm-up worksheet}
\rhead{\small 26 November 2026}
\cfoot{\small\thepage\ / \pageref{LastPage}}
\renewcommand{\headrulewidth}{0.5pt}
\titleformat{\section}{\large\bfseries\color{tgblue}}{\thesection}{0.6em}{}
\titlespacing*{\section}{0pt}{11pt}{5pt}
\setlist[itemize]{leftmargin=*,topsep=3pt,itemsep=2pt,parsep=0pt}
\setlist[enumerate]{leftmargin=*,topsep=3pt,itemsep=3pt,parsep=0pt}

% Blank cell for hand-filling a small matrix: \anscell
\newcommand{\anscell}{\rule{5em}{0.5pt}}

% Ruled writing space: \answerlines{3} prints three blank lines.
\newcount\anslinecount
\newcommand{\answerlines}[1]{%
  \par\vspace{2pt}%
  \anslinecount=0\relax
  \loop
    \ifnum\anslinecount<#1
      \noindent\rule{\textwidth}{0.4pt}\par\vspace{24pt}%
      \advance\anslinecount by 1\relax
  \repeat
  \par\vspace{2pt}}

% Mentors: change to \answerkeytrue to print the answer key at the end.
\newif\ifanswerkey
\answerkeyfalse

\begin{document}
\HackathonHeader{Reconstruct the function}
  {Warm-up worksheet | 09:15--10:15}
  {Participant worksheet}

\section{Introduction}\label{sec:warmup}
\textbf{This one-hour warm-up is ungraded and does not contribute to the final score.}
Given a structural matrix $\mathbf{S}$ and a coefficient matrix $\boldsymbol{\Phi}$,
\textbf{reconstruct the vector-valued function $\mathbf f$ that they encode}:
\[
 \boxed{\quad \mathbf f(\mathbf x)=\boldsymbol{\Phi}\,\mathbf m(\mathbf x)=\mathbf{0},\qquad
 \mathbf f=(f_1,\ldots,f_E)^\mathsf T=\mathbf{0}.\quad\textrm{(implicit system of equations)}\quad}
\]
The inputs defining the function are $\mathbf{S}$ and $\boldsymbol{\Phi}$. The variables
$\mathbf x=(x_1,\ldots,x_N)^\mathsf T$ follow the row order of $\mathbf{S}$;
numerical values of $\mathbf x$ are supplied only when evaluating the function.

\paragraph{Matrix convention.}
Let $N$ be the number of variables, $M$ the number of product terms (monomials), and $E$ the
number of equations. The input dimensions are
\[
 \mathbf x\in\mathbb R^N,\qquad \mathbf{S}\in\mathbb R^{N\times M},\qquad
 \boldsymbol{\Phi}\in\mathbb R^{E\times M}.
\]
\textbf{Rows of $\mathbf{S}$ identify variables; columns identify terms.} Column $j$ of
$\mathbf{S}$ and column $j$ of $\boldsymbol{\Phi}$ refer to the same term. Row $k$ of $\boldsymbol{\Phi}$ supplies
its coefficients in equation $k$. Preserve the supplied variable order.

In TenSyGrid's CPN formulation, $\mathbf{S}$ is binary and the monomials are nonnegative:
\[
 a_{ij}(\mathbf x)=S_{ij}x_i+1-S_{ij},\qquad S_{ij}\in\{0,1\},
 \qquad m_j(\mathbf x)=\prod_{i=1}^{N}a_{ij}(\mathbf x),\qquad
 f_k(\mathbf x)=\sum_{j=1}^{M}\Phi_{kj}m_j(\mathbf x).
\]
An entry $S_{ij}=0$ gives a factor of 1; $S_{ij}=1$ gives $x_i$.
Entries encode factors, \textbf{not powers}, and each column simply selects
which variables appear in a monomial. An all-zero column represents the
constant term 1. Each variable appears at most once in a term, so the
representation is multi-affine in $\mathbf x$. \textbf{Any minus sign belongs to
the coefficient in $\boldsymbol{\Phi}$, never to the monomial.}

\paragraph{Reconstruction and example.}
For each column of $\mathbf{S}$, build its product term. Multiply by the corresponding
coefficients in each row of $\boldsymbol{\Phi}$, then sum to obtain each component $f_k$. Reconstructing $\mathbf f$ does not
require solving $\mathbf f(\mathbf x)=\mathbf 0$.

\textbf{Worked example.} For $\mathbf x=(x_1,x_2)^\mathsf T$, take
\[
 \mathbf{S}=\begin{pmatrix}0&1&0&1\\0&0&1&1\end{pmatrix},\qquad
 \boldsymbol{\Phi}=\begin{pmatrix}-1&1&0&2\\0&0&3&-1\end{pmatrix}.
\]
Then $\mathbf m=(1,x_1,x_2,x_1x_2)^\mathsf T$, giving
\[
 \boxed{\quad f_1(\mathbf x)=-1+x_1+2x_1x_2,\qquad
 f_2(\mathbf x)=3x_2-x_1x_2.\quad}
\]
At $\mathbf x=(1,2)^\mathsf T$, the expected output is $\mathbf f(1,2)=(4,4)^\mathsf T$.

\paragraph{Check your result.}
At 10:15, compare your expressions and evaluator with the reference results and
ask mentors about any differences before moving to grid optimisation. There is
no graded warm-up submission. Check constant terms, zero-valued inputs, signed
coefficients and rectangular systems. Work with the $M$ supplied columns rather than
building all $2^N$ possible monomials.

\newpage
\section{Practice case: from \texorpdfstring{$\mathbf{S}$ and $\boldsymbol{\Phi}$}{S and Phi} to the equations}\label{sec:practice}
The same convention applies, now with the variable vector
\[
 \mathbf x=(\dot x_1,\;x_1,\;u_1,\;y_1,\;z_1)^\mathsf T,
\]
so $N=5$, $M=7$ and $E=2$. The state derivative $\dot x_1$ is a variable in its
own right: the first equation is implicit because it contains the product
$\dot x_1\,y_1$. The supplied structural and coefficient matrices are
\begingroup
\renewcommand{\arraystretch}{1.5}
\setlength{\arraycolsep}{8pt}
\[
 \mathbf{S}=\begin{pmatrix}
  1&0&0&0&0&1&0\\
  0&1&1&1&1&0&0\\
  0&0&1&1&0&0&0\\
  1&0&0&0&0&0&1\\
  0&1&0&1&0&0&0
 \end{pmatrix},\qquad
 \boldsymbol{\Phi}=\begin{pmatrix}
   3& 2&-1&-1& 2& 0& 0\\
   0& 0& 0& 0& 0& 1&-1
 \end{pmatrix}.
\]
\endgroup
Reconstruct the two residual equations $0=f_1(\mathbf x)$ and $0=f_2(\mathbf x)$
encoded by these matrices.

\paragraph{(a) Product terms.} Read each column of $\mathbf{S}$ and write the
corresponding factor product $m_j$:
\par\vspace{3pt}
\begingroup
\renewcommand{\arraystretch}{1.8}
\begin{tabularx}{\textwidth}{@{}c X@{}}
\toprule
\textbf{Column} & \textbf{Product term $m_j(\mathbf x)$}\\
\midrule
$1$ & \\
$2$ & \\
$3$ & \\
$4$ & \\
$5$ & \\
$6$ & \\
$7$ & \\
\bottomrule
\end{tabularx}
\endgroup

\paragraph{(b) Residual equations.} Multiply each term by the coefficient in
the corresponding row of $\boldsymbol{\Phi}$, then simplify:
\vspace{2pt}
\par\noindent $0=f_1(\mathbf x)=$~\hrulefill
\par\vspace{14pt}\noindent\rule{\textwidth}{0.4pt}
\par\vspace{18pt}
\par\noindent $0=f_2(\mathbf x)=$~\hrulefill
\vspace{18pt}

\paragraph{(c) Equilibrium check.} Verify that the residuals vanish at the
equilibrium point $\mathbf x=(2,2,3,2,5)^\mathsf T$:
\vspace{2pt}
\par\noindent $f_1(2,2,3,2,5)=$~\hrulefill
\vspace{14pt}
\par\noindent $f_2(2,2,3,2,5)=$~\hrulefill
\vspace{14pt}

\newpage
\section{Inverse exercise: from equations to \texorpdfstring{$\mathbf{S}$ and $\boldsymbol{\Phi}$}{S and Phi}}\label{sec:inverse}
Now reverse the direction. You are given the two residual equations
\begin{align*}
  0 &= 2 + 3x_1x_2 - 4u_1,\\
  0 &= -1 - x_1x_2 + 5x_2,
\end{align*}
with variable vector $\mathbf x=(x_1,x_2,u_1)^\mathsf T$, so $N=3$ and $E=2$.
Write them as $\mathbf 0=\boldsymbol{\Phi}\,\mathbf m(\mathbf x)$ using exactly the four
factors, in this order,
\[
 \mathbf m(\mathbf x)=\bigl(1,\;x_1x_2,\;u_1,\;x_2\bigr)^\mathsf T,
\]
so $M=4$. Entries of $\mathbf{S}$ are only 0 or 1: a 1 means that the variable appears in
the monomial, an all-zero column gives the constant factor 1, and every minus
sign is placed in $\boldsymbol{\Phi}$.

\paragraph{(a) Structural matrix $\mathbf{S}$.} Rows follow $\mathbf x=(x_1,x_2,u_1)$;
columns follow the prescribed factor order:
\begingroup
\renewcommand{\arraystretch}{2}
\setlength{\arraycolsep}{16pt}
\[
 \mathbf{S}=\begin{pmatrix}
  \anscell&\anscell&\anscell&\anscell\\[18pt]
  \anscell&\anscell&\anscell&\anscell\\[18pt]
  \anscell&\anscell&\anscell&\anscell
 \end{pmatrix}
\]
\endgroup
\vspace{6pt}

\paragraph{(b) Coefficient matrix $\boldsymbol{\Phi}$.} Row $k$ holds the coefficients of
equation $k$; column $j$ refers to the same factor as column $j$ of $\mathbf{S}$:
\begingroup
\renewcommand{\arraystretch}{2}
\setlength{\arraycolsep}{16pt}
\[
 \boldsymbol{\Phi}=\begin{pmatrix}
  \anscell&\anscell&\anscell&\anscell\\[18pt]
  \anscell&\anscell&\anscell&\anscell
 \end{pmatrix}
\]
\endgroup
\vspace{6pt}

\paragraph{(c) Check.} Expand $\boldsymbol{\Phi}\,\mathbf m(\mathbf x)$ and confirm that you
recover both equations:
\answerlines{2}

\vfill
{\par TenSyGrid Hackathon 2026}

\ifanswerkey
\newpage
\section*{Mentor answer key --- do not distribute}
\addcontentsline{toc}{section}{Mentor answer key}
With the column order used above, the factor vector is
\[
 \mathbf m(\mathbf x)=(\dot x_1y_1,\;z_1x_1,\;u_1x_1,\;u_1z_1x_1,\;x_1,\;\dot x_1,\;y_1)^\mathsf T.
\]
Therefore
\[
 \boxed{\quad
 \begin{aligned}
  0 &= 3\dot x_1 y_1 + 2z_1 x_1 - u_1 x_1 - u_1 z_1 x_1 + 2x_1,\\
  0 &= \dot x_1 - y_1.
 \end{aligned}\quad}
\]
At the equilibrium $\mathbf x=(2,2,3,2,5)^\mathsf T$,
$\mathbf m=(4,10,6,30,2,2,2)^\mathsf T$ and
$\boxed{\;f_1=0,\quad f_2=0\;}$: the point satisfies $\dot x_1=y_1=2$ and
$3\dot x_1y_1+x_1(2z_1-u_1-u_1z_1+2)=12-12=0$.
Common pitfalls: confusing $\dot x_1$ with $x_1$, forgetting that an entry
$S_{ij}=0$ contributes a factor 1, putting a minus sign in $\mathbf{S}$ instead of
$\boldsymbol{\Phi}$, and taking coefficients from the wrong row of $\boldsymbol{\Phi}$.

\section*{Inverse exercise --- solution}
With the prescribed column order $(1,x_1x_2,u_1,x_2)$,
\[
 \mathbf{S}=\begin{pmatrix}
  0&1&0&0\\
  0&1&0&1\\
  0&0&1&0
 \end{pmatrix},\qquad
 \boldsymbol{\Phi}=\begin{pmatrix}
  2& 3&-4&0\\
 -1&-1& 0&5
 \end{pmatrix}.
\]
Indeed,
\[
 \boldsymbol{\Phi}\,\mathbf m(\mathbf x)=
 \begin{pmatrix}
  2+3x_1x_2-4u_1\\
  -1-x_1x_2+5x_2
 \end{pmatrix}.
\]
Fixing the order of $\mathbf m$ makes the answer unique; the minus signs are
coefficients in $\boldsymbol{\Phi}$, not part of the monomials.
\fi

\end{document}
